\begin{notechapter}
\chapter{Gated Recurrent Networks: LSTM and GRU}\label{ch:gated-rnn}
\chaptersubtitle{An additive path through time, opened and closed by learned gates}

\begin{chapterintro}
In a simple RNN, the loss at time $s$ reaches the state at time $t$ through a product of $s - t$
Jacobians $\diag(\phi'(a_r))\,W_{hh}$, and such products usually shrink exponentially
(\cref{rnn-bptt:sec:vanishing}). The long short-term memory (LSTM) adds a second state, the
\emph{cell}, which is updated by addition instead of by a matrix product followed by a
squashing function. Along the cell, the Jacobian of one step is a diagonal matrix of
learned ``forget'' values, which the network can keep close to $1$. The gated recurrent unit (GRU)
obtains the same effect with a single state. You will learn
\begin{itemize}[nosep]
  \item the LSTM equations, the role of each gate, and the shapes and parameter count of
        a PyTorch \texttt{nn.LSTM};
  \item the Jacobian of the cell update, the ``constant error carousel'', and what is
        and is not approximated when it is said that LSTMs do not suffer from vanishing
        gradients;
  \item a two-step numerical example, and why the forget-gate bias sets the memory time;
  \item the GRU, and the two different conventions in which it is written;
  \item how to check a hand-written cell against the library.
\end{itemize}
\prerequisites{\cref{ch:rnn-bptt}; the logistic sigmoid (\cref{ch:sigmoid-softmax}).}
\end{chapterintro}

\section{The idea: an additive memory}\label{gated-rnn:sec:idea}

Suppose the state were updated by \emph{adding} something to it, $c_t = c_{t-1} + u_t$, where $u_t$
does not depend on $c_{t-1}$. Then $\partial c_t/\partial c_{t-1} = I$, and a gradient arriving at $c_s$ reaches $c_t$
unchanged, however large $s - t$ is. Hochreiter and Schmidhuber (\emph{Long Short-Term Memory}, 1997), in the paper that introduced
the LSTM, called this identity path the \emph{constant error carousel}. A pure accumulator is useless,
however: it can never forget, and every input is written into it. The LSTM therefore puts
three multiplicative \emph{gates} around it, each a vector of numbers in $(0,1)$ computed from the
current input and the previous output:
\begin{itemize}[nosep]
  \item the \textbf{input gate} decides how much of the new candidate content is written;
  \item the \textbf{forget gate} decides how much of the old cell is kept;
  \item the \textbf{output gate} decides how much of the cell is exposed to the rest of
        the network.
\end{itemize}
The original LSTM of 1997 had input and output gates only; the forget gate was added by
Gers, Schmidhuber and Cummins (\emph{Learning to Forget: Continual Prediction with LSTM},
2000), and the version with forget gate is the one used today.

\section{The LSTM equations}\label{gated-rnn:sec:lstm}

\begin{definition}[LSTM cell]\label{gated-rnn:def:lstm}
Let $x_t\in\R^n$ be the input, and $h_{t-1}\in\R^m$, $c_{t-1}\in\R^m$ the previous output and cell (both $0$ at $t = 0$
by default). With $\sigma$ the logistic sigmoid, applied coordinatewise,
\begin{align*}
  i_t &= \sigma(W_{ii}x_t + W_{hi}h_{t-1} + b_i) &&\text{(input gate)}\\
  f_t &= \sigma(W_{if}x_t + W_{hf}h_{t-1} + b_f) &&\text{(forget gate)}\\
  g_t &= \tanh(W_{ig}x_t + W_{hg}h_{t-1} + b_g) &&\text{(candidate content)}\\
  o_t &= \sigma(W_{io}x_t + W_{ho}h_{t-1} + b_o) &&\text{(output gate)}\\
  c_t &= f_t\odot c_{t-1} + i_t\odot g_t &&\text{(cell)}\\
  h_t &= o_t\odot\tanh(c_t) &&\text{(output, or hidden state)}
\end{align*}
with $W_{i\bullet}\in\R^{m\times n}$, $W_{h\bullet}\in\R^{m\times m}$ and $b_\bullet\in\R^m$.
\end{definition}

The four pre-activations are four affine maps of the same vector $(x_t, h_{t-1})$, so they are
computed by one matrix product with the stacked weights
\[
  \begin{bmatrix}W_{ii} & W_{hi}\\ W_{if} & W_{hf}\\ W_{ig} & W_{hg}\\ W_{io} & W_{ho}\end{bmatrix}\in\R^{4m\times(n+m)},
\]
which is how libraries store them. The cell has $4m(n + m + 1)$ parameters. PyTorch keeps two bias vectors,
\texttt{bias\_ih} and \texttt{bias\_hh}, which are simply added, so \texttt{nn.LSTM} reports
$4m(n + m + 2)$: one of the two is redundant. Its stacking order is $(i, f, g, o)$.

\begin{figure}[ht]
\centering
\begin{tikzpicture}[
  op/.style={draw, circle, inner sep=0pt, minimum size=5.5mm, fill=white, font=\small},
  gate/.style={draw, rounded corners=2pt, minimum width=8mm, minimum height=6mm, font=\scriptsize, fill=bookaccent!15},
  tn/.style={draw, rounded corners=2pt, minimum width=8mm, minimum height=6mm, font=\scriptsize, fill=bookexample!20},
  ar/.style={-{Stealth[length=1.8mm]}, thick},
  lab/.style={font=\scriptsize, inner sep=1pt}]
  \fill[gray!6, rounded corners=6pt] (-0.4,-0.6) rectangle (8.4,3.9);
  % cell line
  \node (cin) at (-1.3,3.2) {$c_{t-1}$};
  \node[op] (mf) at (1.5,3.2) {$\times$};
  \node[op] (pl) at (4.0,3.2) {$+$};
  \node (cout) at (9.5,3.2) {$c_t$};
  \draw[ar, line width=1.4pt, bookpitfall!80!black] (cin) -- (mf);
  \draw[ar, line width=1.4pt, bookpitfall!80!black] (mf) -- (pl);
  \draw[ar, line width=1.4pt, bookpitfall!80!black] (pl) -- (cout);
  % gates
  \node[gate] (f) at (1.5,0.6) {$\sigma$};
  \node[gate] (i) at (3.0,0.6) {$\sigma$};
  \node[tn] (g) at (4.5,0.6) {$\tanh$};
  \node[gate] (o) at (6.0,0.6) {$\sigma$};
  \node[lab, left=1pt of f] {$f_t$};
  \node[lab, left=1pt of i] {$i_t$};
  \node[lab, left=1pt of g] {$g_t$};
  \node[lab, left=1pt of o] {$o_t$};
  \node[op] (mi) at (3.75,1.9) {$\times$};
  \draw[ar] (f) -- (mf);
  \draw[ar] (i) |- (mi); \draw[ar] (g) |- (mi);
  \draw[ar] (mi) -- (pl);
  \node[tn] (tc) at (7.2,2.4) {$\tanh$};
  \draw[ar] (7.2,3.2) -- (tc);
  \node[op] (mo) at (7.2,1.2) {$\times$};
  \draw[ar] (tc) -- (mo); \draw[ar] (o) -- (mo);
  \node (hout) at (9.5,1.2) {$h_t$};
  \draw[ar] (mo) -- (hout);
  % inputs
  \node (hin) at (-1.3,-0.2) {$h_{t-1}$};
  \node (xin) at (0.2,-1.3) {$x_t$};
  \draw[thick] (hin) -- (0.2,-0.2);
  \draw[thick] (xin) -- (0.2,-0.2) -- (6.0,-0.2);
  \foreach \n in {f,i,g,o} \draw[ar] (\n|- 0,-0.2) -- (\n);
  \draw[ar] (8.0,1.2) -- (8.0,-1.0) node[below, lab]{to the next step and to the output layer};
\end{tikzpicture}
\caption{The LSTM cell. The thick line is the cell path: one coordinatewise multiplication by $f_t$
and one addition per step, with no weight matrix and no squashing function. The gates (blue)
and the candidate (green) are computed from $(x_t, h_{t-1})$.}
\label{gated-rnn:fig:lstm}
\end{figure}

\paragraph{Reading the equations.} Each coordinate of the cell is a leaky accumulator:
$c_{t,j} = f_{t,j}\,c_{t-1,j} + i_{t,j}\,g_{t,j}$. With $f_{t,j} = 1$ and $i_{t,j} = 0$ the value is kept exactly; with $f_{t,j} = 0$ and
$i_{t,j} = 1$ it is overwritten by $g_{t,j}\in(-1,1)$. Between the two, it decays geometrically and absorbs
new content. The output $h_t$ is a gated, squashed view of the cell; it is what the next
step's gates see and what the output layer uses. The cell itself never leaves the unit.

\section{The gradient along the cell}\label{gated-rnn:sec:gradient}

In the unrolled graph, $c_{t-1}$ has two children: $c_t$ (through the product $f_t\odot c_{t-1}$) and $h_{t-1}$
(through $\tanh(c_{t-1})$), and $h_{t-1}$ in turn feeds the four gates of step $t$.

\begin{proposition}[Local Jacobian of the cell update]\label{gated-rnn:prop:cell}
With $h_{t-1}$ and $x_t$ held fixed (so that $f_t$, $i_t$, $g_t$ are fixed),
\[
  \frac{\partial c_t}{\partial c_{t-1}} = \diag(f_t) .
\]
The total derivative of $c_t$ with respect to $c_{t-1}$ adds the paths through $h_{t-1}$:
\begin{multline*}
  \frac{dc_t}{dc_{t-1}} = \diag(f_t) + \Bigl(\diag(c_{t-1})\frac{\partial f_t}{\partial h_{t-1}} + \diag(g_t)\frac{\partial i_t}{\partial h_{t-1}}
  + \diag(i_t)\frac{\partial g_t}{\partial h_{t-1}}\Bigr)\\
  \cdot\diag(o_{t-1})\diag\bigl(1 - \tanh^2(c_{t-1})\bigr),
\end{multline*}
where, for instance, $\partial f_t/\partial h_{t-1} = \diag\bigl(f_t\odot(1 - f_t)\bigr)W_{hf}$.
\end{proposition}

\begin{proof}
The first formula: $c_t = f_t\odot c_{t-1} + i_t\odot g_t$ is linear in $c_{t-1}$ with diagonal coefficient $f_t$ when the gates
are fixed. For the second, $f_t$, $i_t$, $g_t$ depend on $c_{t-1}$ only through $h_{t-1} = o_{t-1}\odot\tanh(c_{t-1})$, in which $o_{t-1}$
does not depend on $c_{t-1}$; the product rule on each term of $c_t$ and the chain rule through $h_{t-1}$ give the
bracket.
\end{proof}

The first term is the constant error carousel: along the cell path alone, the loss at
time $s$ reaches $c_t$ through
\[
  \diag(f_s)\diag(f_{s-1})\cdots\diag(f_{t+1}) = \diag\bigl(f_s\odot f_{s-1}\odot\cdots\odot f_{t+1}\bigr),
\]
a product of numbers, coordinate by coordinate, with no matrix $W_{hh}$ and no factor $\phi'$. If the
network keeps $f\approx1$ on a coordinate, this product stays close to $1$ over long distances, and
the gradient carries information about long-range dependencies.

\begin{pitfall}[``LSTMs have no vanishing gradient'']
This statement is too strong, in two ways.
\begin{itemize}[nosep]
  \item Along the cell path, the factor is $\prod f$. It is close to $1$ only on the coordinates and
        time spans where the network has \emph{learned} to keep its forget gates open. On the
        others the contributions vanish, by design: forgetting is the purpose of the gate.
  \item The total Jacobian of \cref{gated-rnn:prop:cell} also contains the paths through $h_{t-1}$, which
        are products of weight matrices and sigmoid derivatives, exactly like those of the
        simple RNN. These paths can vanish or explode; the cell path does not remove them,
        it adds a well-behaved path next to them. The original LSTM went further and
        \emph{truncated} the gradient through the gates, so that only the carousel carried
        errors through time (Hochreiter and Schmidhuber, 1997); modern implementations compute the
        full gradient.
\end{itemize}
The precise statement is: the LSTM \emph{can} propagate gradient information over long
distances along the coordinates where its forget gates are near $1$.
\end{pitfall}

\section{A worked example: memory time}\label{gated-rnn:sec:example}

\begin{example}[A scalar LSTM remembers a pulse]\label{gated-rnn:ex:memory}
Take $m = n = 1$, all recurrent weights $W_{h\bullet} = 0$ (so the gates depend only on $x_t$), input weights
$W_{ii} = W_{ig} = W_{io} = 1$, $W_{if} = 0$, biases $b_i = b_g = b_o = 0$, a forget bias $b_f$, and the input sequence
$x = (1, 0, 0, \dots)$: a pulse followed by silence.

\emph{Step $1$.} $i_1 = \sigma(1) = 0.7311$, $g_1 = \tanh(1) = 0.7616$, $o_1 = \sigma(1) = 0.7311$, and $c_0 = 0$, so
\[
  c_1 = i_1g_1 = 0.5568,\qquad h_1 = o_1\tanh(c_1) = 0.7311\times0.5055 = 0.3696 .
\]
\emph{Steps $t\geq2$.} $x_t = 0$ gives $g_t = \tanh(0) = 0$, so nothing is written, and $f_t = \sigma(b_f)$ is constant: the cell
decays geometrically, $c_t = \sigma(b_f)^{t-1}c_1$, while $h_t = \sigma(0)\tanh(c_t) = 0.5\tanh(c_t)$. With $b_f = 2$, $f = 0.8808$:
$c_2 = 0.4904$, $h_2 = 0.2273$, and after five more steps $c_7 = 0.2600$.

\emph{Memory time.} The half-life of the cell, the number of steps after which half of the
pulse is left, is $\tau_{1/2} = \ln2/(-\ln\sigma(b_f))$:
\begin{center}
\begin{tabular}{c|cccc}
$b_f$ & $0$ & $1$ & $2$ & $3$\\\hline
$f = \sigma(b_f)$ & $0.5$ & $0.7311$ & $0.8808$ & $0.9526$\\
$\tau_{1/2}$ (steps) & $1.00$ & $2.21$ & $5.46$ & $14.27$
\end{tabular}
\end{center}
The same factor $f^{s-t}$ multiplies the gradient along the cell from step $s$ back to step $t$.
\end{example}

The example explains a practical rule. With the usual initialization, all biases are near $0$,
so every forget gate starts near $\sigma(0) = 0.5$: the cell halves at every step, and the
carousel is no better than a simple RNN at the start of training.
Jozefowicz, Zaremba and Sutskever (\emph{An Empirical Exploration of Recurrent Network
Architectures}, 2015) found that initializing the forget bias to $1$ (an idea already in
Gers et al., 2000) substantially improves the LSTM, enough to close its gap with the GRU on
their benchmarks, and they recommend it for every LSTM. PyTorch does not do it by default;
with its stacking order $(i, f, g, o)$ the forget slice of each bias vector is the second
quarter:
\begin{lstlisting}[language=Python]
lstm = torch.nn.LSTM(input_size=n, hidden_size=m)
with torch.no_grad():
    for name, b in lstm.named_parameters():
        if name.startswith("bias_ih"):        # bias_ih and bias_hh are added:
            b[m:2 * m].fill_(1.0)              # set one of them, zero the other
        elif name.startswith("bias_hh"):
            b[m:2 * m].fill_(0.0)
\end{lstlisting}

\section{The gated recurrent unit}\label{gated-rnn:sec:gru}

The GRU (Cho et al., \emph{Learning Phrase Representations using RNN
Encoder--Decoder for Statistical Machine Translation}, 2014) merges cell and output into a single state and the input and
forget gates into a single \emph{update gate}.

\begin{definition}[GRU, PyTorch convention]\label{gated-rnn:def:gru}
\begin{align*}
  r_t &= \sigma(W_{ir}x_t + b_{ir} + W_{hr}h_{t-1} + b_{hr}) &&\text{(reset gate)}\\
  z_t &= \sigma(W_{iz}x_t + b_{iz} + W_{hz}h_{t-1} + b_{hz}) &&\text{(update gate)}\\
  n_t &= \tanh\bigl(W_{in}x_t + b_{in} + r_t\odot(W_{hn}h_{t-1} + b_{hn})\bigr) &&\text{(candidate)}\\
  h_t &= (1 - z_t)\odot n_t + z_t\odot h_{t-1} .
\end{align*}
\end{definition}

The last line is a convex combination, coordinate by coordinate, of the old state and the
candidate. With $h_{t-1}$ held fixed in the gates, $\partial h_t/\partial h_{t-1}$ contains the diagonal term $\diag(z_t)$: $z_t$
plays the role of the LSTM forget gate, and $1 - z_t$ that of its input gate, which are thus tied.
The reset gate decides how much of the old state enters the candidate. The GRU has $3m(n + m + 1)$
parameters ($3m(n + m + 2)$ as counted by PyTorch), three quarters of an LSTM of the same width.

\begin{pitfall}[Two GRUs]
The original paper writes the update as $h_t = z_t\odot h_{t-1} + (1 - z_t)\odot\tilde h_t$ with
$\tilde h_t = \tanh\bigl(Wx_t + U(r_t\odot h_{t-1})\bigr)$: the reset gate is applied to $h_{t-1}$ \emph{before} the matrix $U$.
PyTorch (and cuDNN) apply it \emph{after}, $r_t\odot(W_{hn}h_{t-1} + b_{hn})$. The two are different functions,
since $U(r\odot h)\neq r\odot(Uh)$ in general (\cref{gated-rnn:exr:gru}); the second lets the product $W_{hn}h_{t-1}$ be computed
together with the other two gates in one matrix product. Other texts also swap $z_t$ and $1 - z_t$.
Weights trained in one convention cannot be loaded into the other.
\end{pitfall}

\paragraph{LSTM or GRU?} Neither dominates. Chung et al.\ (\emph{Empirical Evaluation of Gated Recurrent Neural
Networks on Sequence Modeling}, 2014) found the two gated units
clearly better than the simple RNN and comparable to each other on music and speech
modelling. In the large architecture search of Jozefowicz et al., the GRU outperformed
the LSTM with default initialization on all tasks except language modelling, the LSTM with
forget bias $1$ nearly matched it, and no architecture consistently beat both.

\section{Checking a cell against the library}\label{gated-rnn:sec:check}

Library implementations hide the stacking and bias conventions. Writing the cell by hand
and comparing on random inputs removes any doubt; this check passes for the LSTM of
\cref{gated-rnn:def:lstm} and the GRU of \cref{gated-rnn:def:gru}:
\begin{lstlisting}[language=Python]
import torch
torch.set_default_dtype(torch.float64); torch.manual_seed(1)
n, m, T = 3, 4, 5
lstm = torch.nn.LSTM(n, m)                    # input of shape (T, batch, n)
x = torch.randn(T, 1, n)
W, U = lstm.weight_ih_l0, lstm.weight_hh_l0   # shapes (4m, n) and (4m, m)
b = lstm.bias_ih_l0 + lstm.bias_hh_l0
h, c, outs = torch.zeros(m), torch.zeros(m), []
for t in range(T):
    i, f, g, o = (W @ x[t, 0] + U @ h + b).chunk(4)  # stacking order (i, f, g, o)
    i, f, g, o = i.sigmoid(), f.sigmoid(), g.tanh(), o.sigmoid()
    c = f * c + i * g
    h = o * c.tanh()
    outs.append(h)
y, (h_n, c_n) = lstm(x)
print(torch.allclose(torch.stack(outs), y[:, 0]))          # True
print(sum(p.numel() for p in lstm.parameters()))           # 144 = 4*4*(3+4+2)
\end{lstlisting}
The module returns \texttt{y}, the outputs $h_1,\dots,h_T$ of the last layer, and the final pair
$(h_T, c_T)$. By default the time axis comes first, shape $(T, \text{batch}, n)$; \texttt{batch\_first=True}
changes this. With \texttt{num\_layers=2}, the outputs $h_t$ of the first layer are the inputs $x_t$ of the
second.

\begin{chaptersummary}
\begin{itemize}
  \item The LSTM adds a cell $c_t = f_t\odot c_{t-1} + i_t\odot g_t$ to the recurrent state; three sigmoid gates control
        forgetting, writing and reading, and $h_t = o_t\odot\tanh(c_t)$.
  \item With the gates held fixed, $\partial c_t/\partial c_{t-1} = \diag(f_t)$: along the cell, the gradient is multiplied by
        products of forget values, not by products of $\diag(\phi')W_{hh}$. This is the constant error carousel.
  \item The total Jacobian also contains paths through the gates, which behave like those of a
        simple RNN; the LSTM does not abolish vanishing gradients, it adds a path that can avoid
        them where the forget gates are open.
  \item The forget bias sets the memory time, $\tau_{1/2} = \ln2/(-\ln\sigma(b_f))$; initializing it to $1$ is
        recommended.
  \item The GRU interpolates $h_t = (1 - z_t)\odot n_t + z_t\odot h_{t-1}$, with $3/4$ of the LSTM's parameters; its
        original and PyTorch forms differ in where the reset gate is applied.
\end{itemize}
\end{chaptersummary}

\section*{Exercises}
\addcontentsline{toc}{section}{Exercises}

\begin{exercise}\label{gated-rnn:exr:1}
Count the parameters of an LSTM layer and of a GRU layer with $n = 100$ inputs and $m = 256$
units, with one bias vector per gate, then as PyTorch counts them.
\solutionlink{gated-rnn:exr:1}
\end{exercise}

\begin{exercise}\label{gated-rnn:exr:2}
Suppose that $f_t = 1$ and $i_t = 0$ (coordinatewise) for $t = 2,\dots,T$. Show that $c_T = c_1$, and that
$\partial c_T/\partial c_1$ computed along the cell path is the identity. What does $h_T$ depend on?
\solutionlink{gated-rnn:exr:2}
\end{exercise}

\begin{exercise}\label{gated-rnn:exr:3}
In \cref{gated-rnn:ex:memory}, which forget bias gives a half-life of $50$ steps? What is the forget value
$f$, and the factor $f^{50}$ applied to the gradient along the cell over $50$ steps?
\solutionlink{gated-rnn:exr:3}
\end{exercise}

\begin{exercise}\label{gated-rnn:exr:4}
Show that a GRU (PyTorch convention) with $z_t = 0$ and $r_t = 1$ for all $t$ is a simple RNN. Which
weights and activation does it have?
\solutionlink{gated-rnn:exr:4}
\end{exercise}

\begin{exercise}\label{gated-rnn:exr:gru}
Let $U = \begin{bmatrix}1 & 1\\ 0 & 1\end{bmatrix}$, $h = (1, 1)$ and $r = (1, 0)$. Compute $U(r\odot h)$ and $r\odot(Uh)$. Conclude that the original and
the PyTorch GRU are different functions.
\solutionlink{gated-rnn:exr:gru}
\end{exercise}
\end{notechapter}
